\documentclass{situationist-preprint}

\title{Probability in Branching Worlds}
\msid{SCI-09-b}
\msfamily{SCI-09}
\msdiscipline{Philosophy of Science}
\mstype{Companion manuscript}
\mskeywords{Everett interpretation; probability; branching; decision theory; Born rule; Bohmian mechanics}

\begin{document}
\maketitle

\begin{abstract}
The probability problem for the Everett interpretation has two parts: the incoherence problem, how there can be probability at all when every outcome occurs, and the quantitative problem, why the probabilities are the Born weights. SCI-09-a answers the quantitative problem by a decision-theoretic derivation. We argue that the same derivation answers the incoherence problem, because probability is identified by its role in rational decision and in the confirmation of hypotheses by frequencies, and Everettian branch weights play both roles. The comparison with Bohmian mechanics confirms this reading.
\end{abstract}

\section{Two problems}

On the Everett interpretation, a measurement with several possible outcomes produces several quasi-classical branches, each containing an outcome and a successor of the observer. If every outcome occurs, there seems to be nothing for a probability to measure. \citet{Greaves2004} separated this \emph{incoherence problem} from the \emph{quantitative problem}: granting that some measure plays the role of probability, why should it be the squared amplitude? SCI-09-a addresses the second. Its Theorem~5.5 shows that a rational agent who expects to branch assigns an event in an equal-amplitude state a weight equal to its Born weight, Corollary~5.7 extends the result to rational amplitudes, and Proposition~5.8, given continuity, to all amplitudes. It might seem to leave the first problem untouched. We argue that it does not.

\section{Probability by its role}

There is no prior notion of probability, fixed independently of what probability does, against which the Everettian weights could be found wanting. Objective probability is identified by its theoretical role, and the core of that role is twofold \citep{Lewis1980,Wallace2012}. First, rational credence about the outcome of a chancy process should equal the chance, so that chances guide action through expected utility. Second, observed relative frequencies are evidence about chances, so that chance hypotheses are confirmed by statistical data.

\begin{proposition}\label{p:role}
The Born weights of the Everett interpretation play both parts of the role of objective probability.
\end{proposition}

\begin{proof}
The first part is the result of Section~5 of SCI-09-a: by Proposition~5.8 of SCI-09-a, the decision weights of a rational branching agent are the squared amplitudes, so the Born weights guide action through expected utility. The second part is Proposition~5.9 of SCI-09-a: an Everettian agent who weights branches by the Born rule and conditionalizes confirms quantum mechanics by observed frequencies in the same way and to the same degree as an agent who takes the Born rule to be a single-world chance law.
\end{proof}

Uncertainty about which outcome will occur is not part of this role. What makes a single-world agent's credences probabilities is that they are constrained by chances and responsive to frequencies, not that she lacks a fact. The Everettian agent lacks no fact, yet her weights do all the work that probabilities do. The incoherence problem assumes that probability requires ignorance; Proposition~\ref{p:role} shows that the work of probability can be done without it \citep{Deutsch1999,Wallace2012}.

\section{The comparison with Bohmian mechanics}

SCI-09-c argues that Bohmian mechanics contains the full Everettian wave function as a dynamical object and that in quantum equilibrium its particle configuration adds nothing that affects predictions. Bohmian probabilities are probabilities of ignorance: the observer does not know the configuration, and the equilibrium distribution over configurations is the Born distribution.

The comparison shows the reverse. In equilibrium, Bohmian probabilities and Everettian weights coincide branch by branch, and both are fixed by the same amplitudes of the same wave function. What the configuration adds is a mark indicating one branch as actual; it does not alter which weights a rational agent should use or which frequencies confirm the theory. Since both parts of the role of probability are played identically in the two theories, the mark adds nothing to the probabilistic content. This agrees with SCI-09-c and with Lemma~4.4 of SCI-09-a.

\begin{objection}
The decision-theoretic weights describe how an agent should care about her successors, a caring measure, and a caring measure is not a probability.
\end{objection}

\begin{reply}
A caring measure that is fixed by rationality rather than by taste, that coincides with the chances of any rival theory on every prediction, and that is confirmed by frequencies exactly as chances are is a probability by every criterion the role supplies. The only further criterion on offer, ignorance of the outcome, is the assumption the incoherence problem begs.
\end{reply}

\section{Conclusion}

Probability is what plays the role of probability, and in a branching world the Born weights play it.

\begin{thebibliography}{99}\small
\bibitem[Deutsch(1999)]{Deutsch1999} Deutsch, D. (1999). Quantum theory of probability and decisions. \emph{Proceedings of the Royal Society of London A}, 455, 3129--3137.
\bibitem[Greaves(2004)]{Greaves2004} Greaves, H. (2004). Understanding Deutsch's probability in a deterministic multiverse. \emph{Studies in History and Philosophy of Modern Physics}, 35(3).
\bibitem[Lewis(1980)]{Lewis1980} Lewis, D. (1980). A subjectivist's guide to objective chance. In R.~C. Jeffrey (ed.), \emph{Studies in Inductive Logic and Probability}, vol.~II. University of California Press.
\bibitem[Wallace(2012)]{Wallace2012} Wallace, D. (2012). \emph{The Emergent Multiverse: Quantum Theory according to the Everett Interpretation}. Oxford University Press.
\end{thebibliography}

\end{document}
